\overline{\Delta R}_{\mathrm{holdout}}=+0.0235006,
\end{equation}
where positive values favour the static condition. Median holdout delta was $+0.0365401$. The two-sided sign-test value was $p=0.3876953$.

Mean train risk delta was $+0.0069341$, while the median train risk delta was $-0.0033575$. The mean holdout-minus-train generalization gap was $+0.0165665$. Mean promotion count was 3.0 and the mean adaptive/static throughput ratio was $0.98551$.

\begin{figure}[ht]
\centering
\begin{tikzpicture}
\begin{axis}[width=.92\textwidth,height=7cm,xlabel={Replicate},ylabel={$\Delta R=R_{adaptive}-R_{static}$},xmin=.5,xmax=12.5,xtick={1,...,12},grid=major,axis line style={DNCblue},tick label style={font=\scriptsize},label style={font=\small\sffamily}]
\addplot+[mark=*,thick] coordinates {(1,0.005193) (2,-0.010924) (3,0.031937) (4,-0.026776) (5,0.060833) (6,0.051920) (7,0.110019) (8,-0.051079) (9,0.041143) (10,0.050920) (11,0.043128) (12,-0.024308)};
\addplot+[mark=square*,dashed] coordinates {(1,0.049915) (2,-0.015083) (3,0.011120) (4,-0.101112) (5,-0.003333) (6,-0.009349) (7,-0.014609) (8,0.081263) (9,0.058506) (10,-0.003382) (11,-0.009626) (12,0.038898)};
\addplot[black,thin] coordinates {(0.5,0) (12.5,0)};
\legend{Holdout delta,Train delta}
\end{axis}
\end{tikzpicture}
\caption{Per-replicate risk deltas. Positive values indicate lower measured risk for the static anchor; negative values indicate lower measured risk for the adaptive condition.}
\end{figure}

\begin{warningbox}[Benchmark verdict]
\code{NO\_HOLDOUT\_ADVANTAGE\_DEMONSTRATED}. The reference challenge establishes a reproducible zero benchmark for the current selector. External testers can preserve the frozen corpus-generation and scoring contract while replacing or extending attack modules.
\end{warningbox}

\chapter[Generalization-Aware Promotion as the Next Gate]{Generalization-Aware Promotion\\as the Next Gate}
The challenge isolates the next algorithmic requirement: selection pressure should reward attack resistance that persists across independent evaluation folds. A direct extension is a multi-fold promotion predicate. For attack folds $D_1,\ldots,D_k$, define candidate risk $R_j(G)$ on fold $j$ and a generalization penalty
\begin{equation}
\Gamma(G)=\operatorname{Var}\bigl(R_1(G),\ldots,R_k(G)\bigr)+\lambda\max_j\bigl(R_j(G)-R_{\mathrm{parent},j}\bigr)_+.
\end{equation}
A candidate score may then use
\begin{equation}
R_{\mathrm{robust}}(G)=\frac1k\sum_{j=1}^{k}R_j(G)+\beta\Gamma(G),
\end{equation}
with promotion requiring improvement on the robust score and preservation of all frozen invariants. This extension can be benchmarked against the v0.1 reference without changing the frozen AEAD anchor.

\chapter{Security and Validation Register}
\section{Current validated properties}
\begin{longtable}{p{.39\textwidth}p{.18\textwidth}p{.34\textwidth}}
\caption{Evidence register.}\\
\toprule
Property & Status & Evidence\\\midrule
\endfirsthead
\toprule Property & Status & Evidence\\\midrule\endhead
Standard primitive conformance & PASS & RFC/FIPS vectors and differential tests\\
AEAD tamper handling & PASS & mutated-tag fail-closed test\\
DNCrypt v0.2 reference suite & PASS & 21/21 reference tests\\
Native crypto + phase entropy repo layout & PASS & conformance harness\\
Live phase entropy health gate & PASS & 48-frame endogenous runtime receipt\\
36D entropy topology health & PASS & live and validation windows\\
DNCrypt phase entropy integration & PASS & 64/64 unique population, message round-trip, promotion receipt\\
Repository-layout suite after challenge & PASS & 54/54 tests\\
Controlled holdout advantage & NO ADVANTAGE DEMONSTRATED & 4/12 adaptive wins, 8/12 static wins\\
Independent external cryptanalysis & NEXT GATE & frozen challenge package prepared\\
Side-channel evaluation & NEXT GATE & dedicated implementation review required\\
\bottomrule
\end{longtable}

\section{External tester protocol}
An independent test should preserve the following controls:
\begin{enumerate}
\item identical master-key and workload distributions across static and adaptive conditions;
\item identical attack budgets;
\item strict separation of selection data and primary holdout data;
\item exact recording of attack modules, random seeds, platform, compiler/interpreter, and library versions;
\item fail-closed integrity/decrypt checks in every condition;
\item publication of per-replicate outcomes rather than summary values alone.
\end{enumerate}
Attack modules may be strengthened or replaced, provided their budget and observables are symmetric across conditions.

\chapter{Reproducibility}
\section{Repository layout}
The canonical implementation contains the following principal paths:
\begin{lstlisting}
src/dncrypt/
  crypto.py
  crypto_algebra_v01.py
  evolution.py
  entropy_topology_v01.py
  phase_entropy_v01.py
native/
  experimental/
  programs/
protocols/
provenance/
evidence/
tests/
tools/
\end{lstlisting}

\section{Reference benchmark invocation}
From a repository checkout with project dependencies installed, the challenge harness is
\begin{lstlisting}[language=bash]
python tools/cryptanalysis_challenge_v01.py
pytest -q
\end{lstlisting}
The challenge receipt records the BLAKE2b commitment
\begin{center}
\ttfamily\small e08595637366fb182c4b6418d4275f5a2f9415211eb95149256f75dd63924c49
\end{center}
for the reference result set.
